\section{Related Work}

\paragraph{Datasets and task definitions.}
This project builds on ScanNet, a richly annotated RGB-D reconstruction dataset for indoor 3D scene understanding~\citep{dai2017scannet}. ReferIt3D introduced natural referring expressions for fine-grained 3D object identification, including the Nr3D/Sr3D settings used by many later grounding works~\citep{achlioptas2020referit3d}. ScanRefer studies the related problem of localizing language-described objects in ScanNet scenes~\citep{chen2020scanrefer}. These benchmarks make 3D grounding harder than category recognition because the target is often one among several same-class objects and the language may require relation, viewpoint, or anchor evidence.

\paragraph{3D visual grounding.}
Modern 3D grounding methods improve both proposal quality and cross-modal feature interaction. SAT adds 2D semantic assistance to the 3D listener~\citep{yang2021sat}. 3D-SPS avoids a strict detect-then-match pipeline by progressively selecting referred points with language guidance~\citep{luo2022sps}. BUTD-DETR decodes grounded objects from bottom-up objectness and top-down language cues rather than only selecting from fixed proposals~\citep{jain2022butd}. These systems are stronger and broader than the lightweight object-centric pipeline studied here; this report does not claim direct leaderboard comparison against them.

\paragraph{Relation reasoning and anchors.}
Natural referring expressions often name both a target and one or more anchors. InstanceRefer combines instance attributes, instance-to-instance relations, and global localization cues after reducing the candidate set~\citep{yuan2021instancerefer}. 3DVG-Transformer explicitly models proposal relations for relation-enhanced proposal generation and cross-modal disambiguation~\citep{zhao2021_3DVG_Transformer}. EDA decouples text into semantic components and densely aligns language with visual features~\citep{wu2023eda}. These works make clear that relation modeling and dense alignment are established ideas. The current project is narrower: it diagnoses coverage and hard-subset failures, then tests a controlled relation-scoring module. The repository includes anchor-chain parsing utilities, but Phase 4.5 diagnostics show that the conditioned scorer does not yet improve multi-anchor accuracy. For that reason, anchor-chain reasoning is treated as future work rather than a completed contribution.

\paragraph{Conditional computation.}
The method direction is related to mixture-of-experts and conditional computation, where a gate selects or weights different expert functions~\citep{jacobs1991adaptive,shazeer2017outrageously,fedus2022switch}. Here, the experts are pairwise relation scorers rather than generic feed-forward layers. The gate is conditioned on language features, which lets the relation module represent multiple relation modes without requiring explicit viewpoint labels.

\paragraph{Hard negatives and counterfactual training.}
Hard-negative objectives are widely used in visual-semantic retrieval and representation learning~\citep{faghri2018vse}. Counterfactually augmented data can also test whether a model learns the difference that matters for a decision~\citep{kaushik2020learning}. This project implements counterfactual relation learning as an auxiliary pilot objective. The current evidence is promising but small, so it is not used as the main claim.
